\documentclass[11pt,twocolumn]{article}
\usepackage[utf8]{inputenc}
\usepackage[T1]{fontenc}
\usepackage[a4paper,margin=2cm]{geometry}
\usepackage{amsmath,amssymb,amsthm,mathtools}
\usepackage[dvipsnames]{xcolor}
\usepackage{booktabs}
\usepackage{hyperref}
\hypersetup{colorlinks=true,linkcolor=blue!70!black,citecolor=green!50!black}
\usepackage{enumitem}
\setlist{nosep,leftmargin=1.5em}

\theoremstyle{plain}
\newtheorem{theorem}{Theorem}
\newtheorem{proposition}[theorem]{Proposition}
\newtheorem{corollary}[theorem]{Corollary}
\newtheorem{lemma}[theorem]{Lemma}
\theoremstyle{definition}
\newtheorem{definition}[theorem]{Definition}
\newtheorem{example}[theorem]{Example}
\theoremstyle{remark}
\newtheorem{remark}[theorem]{Remark}

\title{\textbf{Cramer's Rule in Twin Spaces:\\Functorial Determinants, Cohomological Obstructions,\\and Twisted Volume Geometry}}
\author{Jocelyn Nemb\'e\\
\small Roots Insights Laboratory, Singapore\\
\small National Institute of Management Sciences, Libreville\\
\small \texttt{jnembe777@gmail.com, jnembe@root-insights.org}}
\date{}

\begin{document}
\maketitle

% ============================================================================
\begin{abstract}
We develop linear algebra inside hierarchies of isomorphic fields (twin spaces) generated by group actions. Given a field $K_0$, a group $G$, and a $G$-family of bijections $\{f_s\}_{s \in G}$ of its underlying set, each $K_s = (K_0, \oplus_s, \otimes_s)$ carries transported operations $x \oplus_s y = f_s^{-1}(f_s(x) + f_s(y))$, and $f_s$ is the structure isomorphism $K_s \to K_0$. We prove five main results:
(1) the twin fields form a groupoid and the matrix algebras $M_n(K_s)$ define a functor from this groupoid to the category of algebras (Theorem~\ref{thm:functor});
(2) the determinant is a \emph{natural transformation} in this groupoid, characterized by alternating $s$-multilinearity and naturality (Theorem~\ref{thm:det-natural});
(3) a generalized Cramer rule holds in every twin field, and its formula is natural --- it commutes with all field morphisms (Theorem~\ref{thm:cramer});
(4) the existence of uniform algebraic rules across the hierarchy is governed by a cohomological obstruction in $H^2(G, K_0^\times)$; Cramer's rule is distinguished by having \emph{trivial} obstruction (Theorem~\ref{thm:obstruction});
(5) when $K_0 = \mathbb{R}$ and the $f_s$ are diffeomorphisms, the determinant measures the \emph{twisted volume} of parallelotopes under the deformed metric $g_s = f_s^*g_0$, with the Jacobian of $f_s$ entering as the pointwise volume density (Theorem~\ref{thm:volume}).

We analyze the algorithmic complexity of solving linear systems: transport to the root field $K_0$, solve in $O(n^3)$, transport back, giving total cost $O(n^2 c_{f_s} + n^3 c_0 + n c_{f_s^{-1}})$, which may be asymptotically cheaper than direct computation in $K_s$. We illustrate with four examples: log-sum-exp (statistical mechanics), $\ell_p$-Minkowski (functional analysis), M\"obius (hyperbolic geometry), and $q$-deformation (quantum groups).
\end{abstract}

\noindent\textbf{Keywords:} twin spaces, determinant, natural transformation, Cramer's rule, cohomological obstruction, twisted volume, groupoid


% ============================================================================
\section{Introduction}
% ============================================================================

Classical linear algebra is built over a fixed field. The determinant measures signed volume; Cramer's rule inverts linear systems; matrix factorizations provide efficient algorithms. But what happens when the field itself varies through a coherent hierarchy of isomorphisms?

The twin spaces framework~\cite{nembe2026vol1,nembe2026vol3,nembe2026vol16} provides such a hierarchy: given an automorphism $f_s$ of a field $K_0$, the transported operations $\oplus_s, \otimes_s$ produce a new field $K_s$ isomorphic to $K_0$ but with different algebraic expressions. This construction arises in information theory (orbit entropy~\cite{nembe2026vol16}), in relativity (Einstein velocity addition as twin addition~\cite{nembe2026vol17}), and in machine learning (log-sum-exp as twin addition with $f_s = \exp$).

The guiding question: \emph{which linear-algebraic identities survive transport through the hierarchy, and which acquire obstructions?}

We show that the answer involves category theory (the determinant is a natural transformation), group cohomology (obstructions live in $H^2(G, K_0^\times)$), and differential geometry (determinants measure twisted volumes). Cramer's rule occupies a distinguished position: it globalizes uniformly because its obstruction class is trivial.


% ============================================================================
\section{Twin Spaces and the Twin Groupoid}
\label{sec:twin-spaces}
% ============================================================================

\begin{definition}[Twin Field]\label{def:twin-field}
Let $(K_0, +, \cdot)$ be a commutative field and $G$ a group acting on the underlying set of $K_0$ by bijections $f_s: K_0 \to K_0$ with $f_e = \mathrm{id}$ and $f_{st} = f_s \circ f_t$ (a homomorphism $G \to \mathrm{Sym}(K_0)$). For each $s \in G$, the \textbf{twin field} $K_s$ is the set $K_0$ equipped with operations:
\begin{align}
x \oplus_s y &= f_s^{-1}(f_s(x) + f_s(y)) \\
x \otimes_s y &= f_s^{-1}(f_s(x) \cdot f_s(y))
\end{align}
with identities $0_s = f_s^{-1}(0)$, $1_s = f_s^{-1}(1)$, and inverse $x^{(-1)}_s = f_s^{-1}(f_s(x)^{-1})$.
\end{definition}

\begin{theorem}[Twin Field Structure]\label{thm:twin-structure}
Each $(K_s, \oplus_s, \otimes_s)$ is a field, and $f_s: (K_s, \oplus_s, \otimes_s) \to (K_0, +, \cdot)$ is a field isomorphism (equivalently $f_s^{-1}: K_0 \to K_s$).
\end{theorem}

\begin{proof}
Every axiom (associativity, commutativity, distributivity, identities, inverses) transfers through $f_s$ by conjugation. For instance, associativity of $\oplus_s$: $(x \oplus_s y) \oplus_s z = f_s^{-1}((f_s(x) + f_s(y)) + f_s(z)) = f_s^{-1}(f_s(x) + (f_s(y) + f_s(z))) = x \oplus_s (y \oplus_s z)$, using associativity of $+$ in $K_0$. The other axioms are analogous. The map $f_s$ intertwines the operations by construction: applying $f_s$ to the defining formulas gives $f_s(x \oplus_s y) = f_s(x) + f_s(y)$ and $f_s(x \otimes_s y) = f_s(x) \cdot f_s(y)$, exhibiting $f_s$ as an isomorphism $K_s \to K_0$ (consistently with $1_s = f_s^{-1}(1)$, i.e.\ $f_s(1_s) = 1$).
\end{proof}

\begin{remark}[Standing hypotheses]\label{rem:hypotheses}
Only the bijectivity of the $f_s$ is needed for the groupoid, the matrix functor (\S\ref{sec:twin-spaces}), the determinant, and Cramer's rule (\S\S\ref{sec:det}--\ref{sec:cramer}): each $K_s$ is a field via transport regardless of how $f_s$ interacts with the operations of $K_0$. Three later results require more. The cohomological obstruction (\S\ref{sec:cohomology}) needs the $f_s$ to act on $K_0^\times$ by automorphisms (multiplicativity), so that $H^2(G, K_0^\times)$ is defined. The twisted-volume geometry (\S\ref{sec:geometry}) needs the $f_s$ to be diffeomorphisms. The descent / twisted-forms results (\S\ref{sec:alg-geom}) need the $f_s$ to be \emph{field} automorphisms of $K_0$, and are therefore non-vacuous only when $\mathrm{Aut}(K_0) \neq \{\mathrm{id}\}$---in particular \emph{not} for $K_0 = \mathbb{R}$, whose only field automorphism is the identity. The headline examples ($\exp$, $x \mapsto x^s$, M\"obius, $q$-deformation) are bijections/diffeomorphisms that are generally \emph{not} field automorphisms, and several live on the semifield $\mathbb{R}_{>0}$; they illustrate the algebraic and geometric results, not the descent results.
\end{remark}

\begin{definition}[Twin Groupoid]\label{def:groupoid}
The \textbf{twin groupoid} $\mathcal{T}(G, K_0)$ has objects $(K_s)_{s \in G}$ and morphisms $f_{t,s} := f_t^{-1} \circ f_s: K_s \to K_t$ for each pair $s, t \in G$ (so that $f_t \circ f_{t,s} = f_s$ as maps to the root $K_0$). Composition: $f_{u,t} \circ f_{t,s} = (f_u^{-1} f_t)(f_t^{-1} f_s) = f_u^{-1} f_s = f_{u,s}$. Every morphism is an isomorphism.
\end{definition}

\begin{definition}[Matrix Functor]\label{def:matrix-functor}
For each $s \in G$, the matrix algebra $M_n(K_s)$ has entrywise operations $\oplus_s, \otimes_s$ and product $[A, B]_s$ with $(c_{ij}) = \bigoplus_{s,k} a_{ik} \otimes_s b_{kj}$. The assignment $K_s \mapsto M_n(K_s)$, $f_{t,s} \mapsto M_n(f_{t,s})$ (entrywise application) defines a functor $M_n: \mathcal{T}(G, K_0) \to \mathrm{Alg}$.
\end{definition}

\begin{theorem}[Functoriality]\label{thm:functor}
$M_n$ is a well-defined functor from the twin groupoid to the category of associative unital algebras. It preserves matrix products: $M_n(f_{t,s})([A, B]_s) = [M_n(f_{t,s})(A), M_n(f_{t,s})(B)]_t$.
\end{theorem}

\begin{proof}
$f_{t,s}$ is a field isomorphism $K_s \to K_t$, so entrywise application preserves the matrix operations. Functoriality: $M_n(f_{u,t}) \circ M_n(f_{t,s}) = M_n(f_{u,t} \circ f_{t,s}) = M_n(f_{u,s})$, and $M_n(\mathrm{id}_{K_s}) = \mathrm{id}_{M_n(K_s)}$.
\end{proof}


% ============================================================================
\section{The Determinant as Natural Transformation}
\label{sec:det}
% ============================================================================

\begin{definition}[Twin Determinant]\label{def:det}
For $A \in M_n(K_s)$, the \textbf{twin determinant} is:
\begin{equation}
\det_s(A) = \bigoplus_{s,\sigma \in S_n} \mathrm{sgn}(\sigma) \otimes_s \bigotimes_{s,i=1}^n a_{i,\sigma(i)}
\end{equation}
where $\mathrm{sgn}(\sigma)$ is embedded in $K_s$ via $1_s$ and $\ominus_s 1_s$.
\end{definition}

\begin{theorem}[Determinant is a Natural Transformation]\label{thm:det-natural}
The family $(\det_s)_{s \in G}$ is a natural transformation $\det: M_n(-) \Rightarrow \mathrm{Id}$ in the twin groupoid:
\begin{enumerate}[label=(\alph*)]
    \item $\det_s$ is alternating and $s$-multilinear in the columns of $A$.
    \item $\det_s(I_n) = 1_s$.
    \item \textbf{Naturality:} for all $s, t \in G$ and $A \in M_n(K_s)$:
    \begin{equation}\label{eq:naturality}
    \boxed{\det_t(M_n(f_{t,s})(A)) = f_{t,s}(\det_s(A))}
    \end{equation}
\end{enumerate}
\end{theorem}

\begin{proof}
Throughout, $f_r: K_r \to K_0$ is the structure isomorphism (Theorem~\ref{thm:twin-structure}), so the intrinsic formula of Definition~\ref{def:det} equals the transported determinant
\[
\det\nolimits_r(A) = f_r^{-1}\bigl(\det\nolimits_0(f_r(A))\bigr), \qquad r \in G,
\]
with $f_r$ applied entrywise and $\det_0$ the classical determinant over $K_0$.

(a,b) These transfer from $K_0$ through $f_s$: alternation and $s$-multilinearity are the $f_s^{-1}$-images of the corresponding properties of $\det_0$, and $\det_s(I_n) = f_s^{-1}(\det_0(I_n)) = f_s^{-1}(1) = 1_s$.

(c) Recall $f_{t,s} = f_t^{-1} \circ f_s$ (Definition~\ref{def:groupoid}), so $f_t \circ f_{t,s} = f_s$ entrywise. For $A \in M_n(K_s)$,
\begin{align}
\det\nolimits_t\bigl(M_n(f_{t,s})(A)\bigr)
&= f_t^{-1}\Bigl(\det\nolimits_0\bigl(f_t(f_{t,s}(A))\bigr)\Bigr) \notag\\
&= f_t^{-1}\bigl(\det\nolimits_0(f_s(A))\bigr),
\end{align}
while
\begin{align}
f_{t,s}\bigl(\det\nolimits_s(A)\bigr)
&= f_t^{-1}\Bigl(f_s\bigl(f_s^{-1}(\det\nolimits_0(f_s(A)))\bigr)\Bigr) \notag\\
&= f_t^{-1}\bigl(\det\nolimits_0(f_s(A))\bigr).
\end{align}
The two expressions agree, so the naturality square commutes.
\end{proof}


% ============================================================================
\section{Generalized Cramer's Rule}
\label{sec:cramer}
% ============================================================================

\begin{theorem}[Generalized Cramer Rule]\label{thm:cramer}
Let $A \in M_n(K_s)$ with $\det_s(A) \neq 0_s$. The linear system $[A, x]_s = b$ (where all operations use $\oplus_s, \otimes_s$) has the unique solution:
\begin{equation}
\boxed{x_j = \det_s(B_j) \oslash_s \det_s(A)}
\end{equation}
where $B_j$ is $A$ with the $j$-th column replaced by $b$, and $\oslash_s$ is twin division.

Moreover, this formula is \textbf{natural}: applying $f_{t,s}$ to both sides gives the Cramer solution for the transported system in $K_t$.
\end{theorem}

\begin{proof}
Transport to $K_0$: the system $[A, x]_s = b$ becomes $[\hat{A}, \hat{x}]_0 = \hat{b}$ where $\hat{A} = f_s(A)$, $\hat{b} = f_s(b)$. Classical Cramer: $\hat{x}_j = \det_0(\hat{B}_j) / \det_0(\hat{A})$. Transport back: $x_j = f_s^{-1}(\hat{x}_j) = f_s^{-1}(\det_0(\hat{B}_j) / \det_0(\hat{A})) = \det_s(B_j) \oslash_s \det_s(A)$.

Naturality: applying $f_{t,s}$ and using naturality of $\det$ (\ref{eq:naturality}):
$f_{t,s}(x_j) = f_{t,s}(\det_s(B_j)) \oslash_t f_{t,s}(\det_s(A)) = \det_t(f_{t,s}(B_j)) \oslash_t \det_t(f_{t,s}(A))$, which is the Cramer solution in $K_t$.
\end{proof}


% ============================================================================
\section{Cohomological Obstruction}
\label{sec:cohomology}
% ============================================================================

Not all algebraic identities globalize as cleanly as Cramer's rule. We formalize the obstruction. Throughout this section (cf.\ Remark~\ref{rem:hypotheses}) we assume the $f_s$ are multiplicative on $K_0^\times$, so that $s \cdot a := f_s(a)$ is an action of $G$ on the abelian group $K_0^\times$ by automorphisms and the group cohomology $H^2(G, K_0^\times)$ is defined.

\begin{definition}[Uniform Rule]\label{def:uniform}
A family of maps $(\Phi_s: \mathcal{L}_s \to K_s^n)_{s \in G}$ (where $\mathcal{L}_s$ is the set of invertible linear systems over $K_s$) is \textbf{uniform} if for all $s, t \in G$:
$f_{t,s}(\Phi_s(A, b)) = \Phi_t(f_{t,s}(A), f_{t,s}(b))$.
\end{definition}

\begin{definition}[Obstruction Cocycle]\label{def:cocycle}
Given local rules $\phi_s$ for a property $P$, transport from $K_s$ to $K_t$ produces a discrepancy factor $c_P(t, s) \in K_0^\times$ satisfying the cocycle condition:
$c_P(s,t) \cdot f_s(c_P(t,u)) = c_P(s,tu) \cdot c_P(st,u)$.
The class $[c_P] \in H^2(G, K_0^\times)$ is the \textbf{obstruction class} of $P$.
\end{definition}

\begin{theorem}[Cohomological Obstruction]\label{thm:obstruction}
There exists a uniform rule for $P$ across all twin spaces if and only if $[c_P] = 0$ in $H^2(G, K_0^\times)$.
\end{theorem}

\begin{proof}
$(\Rightarrow)$ If $(\Phi_s)_s$ is uniform: transport produces no discrepancy, so $c_P \equiv 1$ and $[c_P] = 0$.

$(\Leftarrow)$ If $[c_P] = 0$: there exists $b: G \to K_0^\times$ with $c_P(s,t) = b(st) b(s)^{-1} f_s(b(t)^{-1})$. The modified rules $\Phi_s = b(s)^{-1} \cdot \phi_s$ are uniform (the discrepancy is absorbed by $b$).
\end{proof}

\begin{theorem}[Cramer Has Trivial Obstruction]\label{thm:cramer-trivial}
For Cramer's rule, $[c_{\mathrm{Cr}}] = 0$: the obstruction is trivial.
\end{theorem}

\begin{proof}
The Cramer formula $x_j = \det_s(B_j) \oslash_s \det_s(A)$ involves only the determinant and field division. Both are natural (Thm.~\ref{thm:det-natural}), so transport produces no discrepancy: $c_{\mathrm{Cr}}(t,s) = 1$ for all $t, s$.
\end{proof}

\begin{example}[Non-trivial Obstruction]
Consider the Cholesky decomposition $A = L L^T$ in $K_s$. The decomposition requires choosing square roots ($\sqrt{\cdot}_s$ in $K_s$), which depend on the sign conventions of $\oplus_s$. Different twin fields may require different branch choices, producing a non-trivial cocycle $c_{\mathrm{Chol}}(t,s) = \pm 1 \in K_0^\times$. If $G \cong \mathbb{Z}/2\mathbb{Z}$ and $H^2(\mathbb{Z}/2, K_0^\times) \neq 0$, the Cholesky decomposition does not globalize.
\end{example}


% ============================================================================
\begin{example}[A genuinely non-trivial obstruction: $\mathrm{Br}(\mathbb{R})$]\label{ex:brauer-z2}
Take $K_0=\mathbb{C}$ and $G=\mathbb{Z}/2$ acting by complex conjugation $f_\sigma(z)=\bar z$ --- a genuine field
automorphism, so the hypotheses of \S\ref{sec:cohomology} hold. For a cyclic group, $H^2$ is the Tate quotient
\[
H^2(\mathbb{Z}/2,\mathbb{C}^\times)=\frac{(\mathbb{C}^\times)^{G}}{N(\mathbb{C}^\times)}=\frac{\mathbb{R}^\times}{\mathbb{R}_{>0}}\cong\mathbb{Z}/2,
\]
since the fixed subgroup is $(\mathbb{C}^\times)^\sigma=\mathbb{R}^\times$ and the norm $N(z)=z\bar z=|z|^2$ has image
$\mathbb{R}_{>0}$. This is exactly the Brauer group $\mathrm{Br}(\mathbb{R})\cong\mathbb{Z}/2$, whose nontrivial class
is Hamilton's quaternion algebra. A rule whose obstruction cocycle (Definition~\ref{def:cocycle}) represents this
class therefore \emph{cannot} be globalized across the two twin fields $\mathbb{C}$ and $\overline{\mathbb{C}}$,
in sharp contrast to Cramer's rule, whose class is trivial (Theorem~\ref{thm:cramer-trivial}). The obstruction
theory of \S\ref{sec:cohomology} thus literally detects the quaternions, realizing concretely the Brauer-group
connection of \S\ref{sec:alg-geom}.
\end{example}

% ============================================================================
\section{Twisted Volume Geometry}
\label{sec:geometry}
% ============================================================================

When $K_0 = \mathbb{R}$ and each $f_s$ is a $C^\infty$-diffeomorphism:

\begin{definition}[Twisted Metric]\label{def:twisted-metric}
Let each $f_s$ be a $C^\infty$-diffeomorphism, acting componentwise on $K_s^n = \mathbb{R}^n$. The \textbf{twisted metric} is the pullback $g_s := f_s^* g_0$, i.e.\ $g_s(v, w) = g_0(Df_s\, v, Df_s\, w)$ at each point. When $f_s$ is linear this reduces to $g_s(v, w) = g_0(f_s(v), f_s(w))$.
\end{definition}

\begin{theorem}[Determinant and Twisted Volume]\label{thm:volume}
Let each $f_s$ be a $C^1$-diffeomorphism acting componentwise, and $A \in M_n(K_s)$ with columns $v_1, \ldots, v_n$. Then
\begin{equation}
\bigl|f_s(\det_s(A))\bigr| = \bigl|\det\nolimits_0(f_s(A))\bigr| = \mathrm{vol}_0\bigl(P(f_s(v_1), \ldots, f_s(v_n))\bigr),
\end{equation}
the Euclidean volume of the parallelotope spanned by the componentwise-transported edge vectors $f_s(v_j)$. Under the twisted metric $g_s = f_s^* g_0$ this equals the $g_s$-volume of $P(v_1, \ldots, v_n)$ when $f_s$ is linear; for nonlinear $f_s$ it is the infinitesimal (tangent-space) volume, the finite volume $\mathrm{vol}_{g_s}(P) = \int_P |\det Df_s|\, d\mathrm{vol}_0$ acquiring the pointwise Jacobian factor.
\end{theorem}

\begin{proof}
By transport (Theorem~\ref{thm:det-natural}), $\det_s(A) = f_s^{-1}(\det_0(f_s(A)))$, hence $f_s(\det_s(A)) = \det_0(f_s(A))$. The matrix $f_s(A)$ has columns $f_s(v_1), \ldots, f_s(v_n)$ (entrywise $f_s$), so $|\det_0(f_s(A))|$ is the Euclidean volume of the parallelotope they span. The metric statement is the change-of-variables formula for $g_s = f_s^* g_0$: $\mathrm{vol}_{g_s}(P) = \int_P \sqrt{\det g_s} = \int_P |\det Df_s|$, which for linear $f_s$ is $|\det f_s|$ times $\mathrm{vol}_0(P)$.
\end{proof}

\begin{proposition}[Volume Density, not Curvature]\label{prop:density}
Since $g_s = f_s^* g_0$ is the pullback of $g_0$ by the diffeomorphism $f_s$, the map $f_s$ is an isometry $(\mathbb{R}^n, g_s) \to (\mathbb{R}^n, g_0)$. Consequently curvature is \emph{preserved}, $R_s = R_0 \circ f_s$; in particular $g_s$ is flat whenever $g_0$ is. The quantity that varies pointwise is the \emph{volume density}
\[
\sqrt{\det g_s} = |\det Df_s| = \prod_{i=1}^n |f_s'(x_i)|,
\]
the Jacobian of the transport---this density, not a curvature, is the local ``cost'' of computing in twin coordinates.
\end{proposition}

\begin{remark}[Connection to Vol.~17]
The Jacobian volume density $|\det Df_s|$ is the linear-algebra counterpart of the transport factor in the twin formulation of general relativity~\cite{nembe2026vol17}: the ``cost'' of working in twin coordinates enters as a change of measure. It is a density, not a curvature---pullback by a diffeomorphism leaves the intrinsic curvature unchanged.
\end{remark}


% ============================================================================
\section{Algebraic Geometry: Twin Spaces as Twisted Forms}
\label{sec:alg-geom}
% ============================================================================

\begin{theorem}[Twin $\mathrm{GL}_n$ as Twisted Forms]\label{thm:twisted-forms}
The twin hierarchy induces a family of twisted forms $\{\mathrm{GL}_{n,s}\}_{s \in G}$ of the group scheme $\mathrm{GL}_n$ over $K_0$, obtained by base change along $f_s$:
$\mathrm{GL}_{n,s} = \mathrm{GL}_n \times_{\mathrm{Spec}(K_0), f_s} \mathrm{Spec}(K_0)$.

The transition maps $\varphi_{t,s}: \mathrm{GL}_{n,s} \to \mathrm{GL}_{n,t}$ satisfy the cocycle condition $\varphi_{u,t} \circ \varphi_{t,s} = \varphi_{u,s}$ and form descent data in the sense of Grothendieck~\cite{grothendieck1960}.
\end{theorem}

\begin{proof}
Each $f_s$ is a $K_0$-automorphism, inducing base change. The cocycle condition follows from $f_u \circ f_t^{-1} \circ f_t \circ f_s^{-1} = f_u \circ f_s^{-1}$. The determinant extends to a morphism $\det_s: \mathrm{GL}_{n,s} \to \mathbb{G}_{m,s}$ with exact sequence $1 \to \mathrm{SL}_{n,s} \to \mathrm{GL}_{n,s} \xrightarrow{\det_s} \mathbb{G}_{m,s} \to 1$.
\end{proof}

\begin{remark}[Hypothesis for descent]
This is the one place where genuine \emph{field} automorphisms are required: base change along $f_s$ presupposes $f_s \in \mathrm{Aut}(K_0)$ as a field, and the descent datum is Galois descent for the $G$-action (Remark~\ref{rem:hypotheses}). For $K_0 = \mathbb{R}$ the construction is trivial ($\mathrm{Aut}(\mathbb{R}) = \{\mathrm{id}\}$); the natural settings are $K_0 = \mathbb{C}$, finite fields, or function fields, where the Galois group is nontrivial.
\end{remark}

\begin{remark}[Brauer Group Connection]
The twisted forms $\mathrm{GL}_{n,s}$ correspond to classes in $H^1(G, \mathrm{Aut}(\mathrm{GL}_n))$. The same $H^2(G, K_0^\times)$ that controls uniform rules also controls the Brauer group $\mathrm{Br}(K_0)$, connecting the cohomological obstruction (\S\ref{sec:cohomology}) to the geometry of twisted group schemes.
\end{remark}


% ============================================================================
\section{Algorithmic Complexity}
\label{sec:complexity}
% ============================================================================

\begin{definition}[Cost Model]
For twin field $K_s$: $c_{\oplus_s}, c_{\otimes_s}$ = cost of one twin operation; $c_{f_s}, c_{f_s^{-1}}$ = cost of evaluating the automorphism; $c_0$ = cost of one operation in $K_0$.
\end{definition}

\begin{theorem}[Complexity Comparison]\label{thm:complexity}
Solving $[A, x]_s = b$ with $A \in M_n(K_s)$:
\begin{enumerate}[label=(\alph*)]
    \item \textbf{Direct in $K_s$} (Gaussian elimination): $C_{\mathrm{direct}} = O(n^3(c_{\oplus_s} + c_{\otimes_s}))$.
    \item \textbf{Via transport to $K_0$}: $C_{\mathrm{transport}} = O(n^2 c_{f_s} + n^3 c_0 + n c_{f_s^{-1}})$.
\end{enumerate}
Transport is cheaper when $n^3(c_{\oplus_s} + c_{\otimes_s}) > n^2 c_{f_s} + n^3 c_0$, i.e., when twin operations are more expensive than transfer + classical operations.
\end{theorem}

\begin{proof}
(a) Gaussian elimination requires $O(n^3)$ additions and multiplications, each costing $c_{\oplus_s}$ or $c_{\otimes_s}$.
(b) Transport: $n^2 + n$ evaluations of $f_s$ (cost $O(n^2 c_{f_s})$), then $O(n^3)$ classical operations (cost $O(n^3 c_0)$), then $n$ evaluations of $f_s^{-1}$ (cost $O(n c_{f_s^{-1}})$).
\end{proof}

\begin{theorem}[Optimal Resolution Space]\label{thm:optimal-space}
For a system posed in $K_s$, the optimal resolution space is $t^* = \arg\min_{t \in G}[n^2 c_{f_{t,s}} + C_{\mathrm{solve}}(t, n) + n c_{f_{s,t}}]$.
\end{theorem}

\begin{remark}[Connection to Vol.~16]
This is the linear-algebra analog of the optimal group selection problem in twin information theory~\cite{nembe2026vol16,nembe2026a2}: choosing the ``right'' twin space minimizes cost, whether the cost is entropy (information theory) or operations (linear algebra).
\end{remark}


% ============================================================================
\section{Examples}
\label{sec:examples}
% ============================================================================

The following examples illustrate the transport mechanism. Several use the semifield $\mathbb{R}_{>0}$ and bijections that are not field automorphisms; per Remark~\ref{rem:hypotheses} they exhibit the algebraic (determinant, Cramer) and geometric results rather than the Galois-descent results of \S\ref{sec:alg-geom}.

\begin{example}[Log-Sum-Exp (Statistical Mechanics)]\label{ex:logsumexp}
$K_0 = \mathbb{R}_{>0}$, $f_s(x) = e^{sx}$. Then $x \oplus_s y = \frac{1}{s}\log(e^{sx} + e^{sy})$ (log-sum-exp) and $x \otimes_s y = x + y$. The determinant $\det_s(A) = \frac{1}{s}\log(\sum_\sigma \mathrm{sgn}(\sigma) \exp(s \sum_i a_{i,\sigma(i)}))$ appears in partition functions and softmax operators in machine learning.
\end{example}

\begin{example}[Minkowski $\ell_p$ Geometry]\label{ex:minkowski}
$K_0 = \mathbb{R}_{>0}$, $f_s(x) = x^s$. Then $x \oplus_s y = (x^s + y^s)^{1/s}$ ($\ell_s$-addition) and $x \otimes_s y = xy$. For $s = 2$: Euclidean addition $(x^2 + y^2)^{1/2}$. The determinant measures $\ell_s$-volume.
\end{example}

\begin{example}[M\"obius (Hyperbolic Geometry)]\label{ex:mobius}
$f_s(z) = (z + s)/(1 - sz)$ on the unit disk. Operations are rational functions; the determinant encodes oriented hyperbolic area. Cramer's rule in this geometry solves hyperbolic barycentric systems.
\end{example}

\begin{example}[$q$-Deformation (Quantum Groups)]\label{ex:q-deform}
$f_q(x) = (q^x - q^{-x})/(q - q^{-1})$ for $q > 1$. The operations $\oplus_q, \otimes_q$ are $q$-analogs. The determinant $\det_q$ relates to the $q$-determinant in representation theory of quantum groups~\cite{manin1988}.
\end{example}


% ============================================================================
\section{Conclusion}
\label{sec:conclusion}
% ============================================================================

We have shown that linear algebra extends coherently to twin spaces, with the determinant as a natural transformation, Cramer's rule as a globally uniform identity (trivial obstruction), and volumes twisted by the Jacobian of the field automorphism. The cohomological obstruction distinguishes which identities survive transport (Cramer, determinant product formula) from those that may fail (matrix factorizations, spectral decompositions).

\textbf{Open problems:} (1) Galois theory for twin field extensions. (2) Twin Hilbert spaces with twisted spectra. (3) Numerical algorithms exploiting optimal twin geometry for large-scale systems. (4) Connections to gauge theory: the twin groupoid as a gauge groupoid, with the determinant as a gauge-invariant observable.


% ============================================================================
\begin{thebibliography}{15}

\bibitem{nembe2026vol1} J.~Nemb\'e, ``Twin Spaces: Foundations,'' \emph{ROOTS-INSIGHTS}, Vol.~1, 2026.

\bibitem{nembe2026vol3} J.~Nemb\'e, ``The Transfer Theorem,'' \emph{ROOTS-INSIGHTS}, Vol.~3, 2026.

\bibitem{nembe2026vol16} J.~Nemb\'e, ``Twin Information Theory,'' \emph{ROOTS-INSIGHTS}, Vol.~16, 2026.

\bibitem{nembe2026vol17} J.~Nemb\'e, ``Twin General Relativity,'' \emph{ROOTS-INSIGHTS}, Vol.~17, 2026.

\bibitem{nembe2026a2} J.~Nemb\'e, ``Optimal Group Selection for Communication,'' 2026.

\bibitem{grothendieck1960} A.~Grothendieck, ``Technique de descente et th\'eor\`emes d'existence en g\'eom\'etrie alg\'ebrique,'' \emph{S\'em.\ Bourbaki}, 1959--1960.

\bibitem{serre1979} J.-P.~Serre, \emph{Local Fields}, Springer, 1979.

\bibitem{milne2017} J.S.~Milne, \emph{Algebraic Geometry}, online notes, 2017.

\bibitem{manin1988} Yu.I.~Manin, \emph{Quantum Groups and Non-commutative Geometry}, CRM, 1988.

\bibitem{lang2002} S.~Lang, \emph{Algebra}, 3rd ed., Springer, 2002.

\bibitem{neukirch1999} J.~Neukirch, \emph{Algebraic Number Theory}, Springer, 1999.

\bibitem{docarmo1992} M.P.~do~Carmo, \emph{Riemannian Geometry}, Birkh\"auser, 1992.

\end{thebibliography}

\end{document}
